\section*{Hoofdstuk \ref{ch:g1:landscapes-around-us} --- Landschappen om ons heen}

\begin{solution}{exo:g1:landscapes-around-us:1}
\omterm{def:g1:landscapes-around-us:landform}{Berg}, \omterm{def:g1:landscapes-around-us:landform}{heuvel}, \omterm{def:g1:landscapes-around-us:landform}{dal}, \omterm{def:g1:landscapes-around-us:landform}{vlakte}.
\end{solution}

\begin{solution}{exo:g1:landscapes-around-us:2}
Een \omterm{def:g1:landscapes-around-us:landform}{berg}: die is hoger en steiler dan een \omterm{def:g1:landscapes-around-us:landform}{heuvel}.
\end{solution}

\begin{solution}{exo:g1:landscapes-around-us:3}
Natuurlijk: de \omterm{def:g1:landscapes-around-us:water}{rivier}, het bos, het strand. Door mensen gemaakt: de weg, de brug, het
huis.
\end{solution}

\begin{solution}{exo:g1:landscapes-around-us:4}
Waar het land de zee raakt.
\end{solution}

\begin{solution}{exo:g1:landscapes-around-us:5}
Bijvoorbeeld: natuurlijk zijn de \omterm{def:g1:landscapes-around-us:landform}{bergen}, het bos op de hellingen of de
beek; door mensen gemaakt zijn een huis of de kerk van het dorp.
\end{solution}

\begin{solution}{exo:g1:landscapes-around-us:6}
De \omterm{def:g1:landscapes-around-us:landform}{berg}, met sneeuw op zijn top. De \omterm{def:g1:landscapes-around-us:water}{rivier} stroomt onderin het \omterm{def:g1:landscapes-around-us:landform}{dal}, tussen
de \omterm{def:g1:landscapes-around-us:landform}{bergen} en de \omterm{def:g1:landscapes-around-us:landform}{heuvel}.
\end{solution}

\begin{solution}{exo:g1:landscapes-around-us:7}
Een \omterm{def:g1:landscapes-around-us:water}{rivier} stroomt van hoger land (\omterm{def:g1:landscapes-around-us:landform}{heuvels} en \omterm{def:g1:landscapes-around-us:landform}{bergen}) naar beneden, tot in
een \omterm{def:g1:landscapes-around-us:water}{meer} of in de zee. Water stroomt altijd naar beneden, dus het kan niet
omhoog naar de top van een \omterm{def:g1:landscapes-around-us:landform}{heuvel} stromen.
\end{solution}

\begin{solution}{exo:g1:landscapes-around-us:8}
De boom is dichtbij, op de voorgrond; de \omterm{def:g1:landscapes-around-us:landform}{heuvel} is ver weg, op de
achtergrond. Dingen die ver weg zijn, lijken kleiner dan dingen die
dichtbij zijn.
\end{solution}

\begin{solution}{exo:g1:landscapes-around-us:9}
Voorgrond: hoog gras en een paar struiken onderaan het plaatje.
Middengrond: de \omterm{def:g1:landscapes-around-us:water}{rivier} die over de \omterm{def:g1:landscapes-around-us:landform}{vlakte} kronkelt, met bomen langs de
oever, akkers, een boerderij en een weg.
Achtergrond: het vlakke land dat doorloopt tot aan de horizon, en de hemel.
\end{solution}

\begin{solution}{exo:g1:landscapes-around-us:10}
Omdat mensen het gemaakt hebben: ze hebben het land vrijgemaakt, geploegd
en de tarwe gezaaid. Zonder mensen zou daar gras, struikgewas of bos
groeien, geen tarweveld.
\end{solution}

\begin{solution}{exo:g1:landscapes-around-us:11}
Niet veranderd: de \omterm{def:g1:landscapes-around-us:landform}{heuvel} en de \omterm{def:g1:landscapes-around-us:water}{rivier}. Toegevoegd door mensen: de akkers,
meer huizen, de weg en de brug (en het grootste deel van het bos is gekapt).
\end{solution}
